% Problem record op_72937f30974953ba. This TeX file is derived from
% database/problems_json/op_72937f30974953ba.json by scripts/sync-tex.mjs; edit the JSON record
% and rerun the script rather than editing this file.
\section{Conditional-correlation decay in amplitude-damped random circuits}
\label{sec:op_72937f30974953ba}

\paragraph{Problem.}
Do computational-basis outputs of one-dimensional Haar-random circuits with fixed amplitude damping have conditional mutual information that decays exponentially with separation, uniformly in circuit depth?

Consider an $n$-qubit, depth-$d$ nearest-neighbor Haar-random brickwork circuit on a line, starting from $|0^n\rangle$. After each layer, apply independently to every qubit the amplitude-damping channel of fixed strength $\gamma\in(0,1)$, with Kraus operators

\begin{equation}
K_0:=|0\rangle\langle0|+\sqrt{1-\gamma}|1\rangle\langle1|,
\qquad
K_1:=\sqrt\gamma|0\rangle\langle1|.
\label{eq:7293-1}
\end{equation}

Let $P_U$ be the computational-basis output distribution for the collection $U$ of random gates. For disjoint subsets $A,B,C\subseteq\{1,\ldots,n\}$, define the classical conditional mutual information

\begin{equation}
I_{P_U}(A:C|B):=H_{P_U}(AB)+H_{P_U}(BC)-H_{P_U}(B)-H_{P_U}(ABC),
\label{eq:7293-2}
\end{equation}

where $H_{P_U}$ is Shannon entropy in bits. Let $r(A,C)$ be the minimum chain distance between nonempty $A$ and $C$.

For every fixed $\gamma\in(0,1)$, are there constants $a,b,c>0$, independent of $n,d,A,B,C$, such that

\begin{equation}
\mathbb E_U I_{P_U}(A:C|B)
\leq a n^b e^{-c r(A,C)}
\label{eq:7293-3}
\end{equation}

for every depth and every such choice of disjoint subsets?

The displayed definitions, constraints, and target bounds are recorded in Eqs.~\eqref{eq:7293-1}, \eqref{eq:7293-2}, \eqref{eq:7293-3}.

\subsection*{Status}

Unsolved

\subsection*{Source}

This precise formulation is editor wording based on the unresolved direction and limitations documented in the cited primary literature \sourcecite{ref:7293-lee25}{Lee25}\sourcecite{ref:7293-mele26}{Mele26}\sourcecite{ref:7293-shravan26}{Shravan26}; it is not presented as a verbatim conjecture of those authors.

\subsection*{Progress}

\begin{itemize}
  \item Lee and coauthors identify this average approximate-Markov condition as sufficient for classical sampling. In one dimension, their Theorem 3 gives runtime $\operatorname{poly}(n,1/\varepsilon,1/\delta)$ and output error $\|P_U-Q_U\|_1\leq\varepsilon$, except on a fraction $\delta$ of circuits, provided the condition holds uniformly over depth. Their amplitude-damping evidence is numerical, not a general proof. \sourcecite{ref:7293-lee25}{Lee25}

  \item Mele and coauthors prove effective-depth bounds for expectation values under nonunital noise. For a bounded observable $O$, the effect of discarding all but the last $m$ noisy layers is bounded on average by
  \begin{equation}
O\!\left(\|O\|_\infty e^{-\alpha_\gamma m}\right),
  \qquad \alpha_\gamma>0.
\label{eq:7293-4}
\end{equation}
  Their 2026 publication explicitly distinguishes these results from the still-open general sampling problem. \sourcecite{ref:7293-mele26}{Mele26}

The displayed definitions, constraints, and target bounds are recorded in Eqs.~\eqref{eq:7293-4}.

  \item The April 2026 IQP simulation paper explicitly describes the nonunital conditional-mutual-information evidence as numerical. Its polynomial-time result for amplitude-damped IQP circuits at $d=\Omega(\log n)$ uses a different, diagonal-gate structure and does not prove the Haar-circuit inequality above. No later proof of this inequality was located. \sourcecite{ref:7293-shravan26}{Shravan26}
\end{itemize}

\subsection*{References}

\begin{enumerate}[label={},leftmargin=4.8em,labelsep=0.5em]
  \footnotesize
  \item[\textup{[Lee25]}]\label{ref:7293-lee25}
  S.-u. Lee, S. Ghosh, C. Oh, K. Noh, B. Fefferman, and L. Jiang, "Classical simulation of noisy random circuits from exponential decay of correlation," arXiv preprint (2025), version 1, 7 October 2025. \href{https://arxiv.org/abs/2510.06328}{arXiv:2510.06328}.

  \item[\textup{[Mele26]}]\label{ref:7293-mele26}
  A. A. Mele, A. Angrisani, S. Ghosh, S. Khatri, J. Eisert, D. Stilck França, and Y. Quek, "Noise-induced shallow circuits and the absence of barren plateaus," \emph{Nature Physics} 22, 751–756 (2026). \href{https://doi.org/10.1038/s41567-026-03245-z}{doi:10.1038/s41567-026-03245-z}; \href{https://arxiv.org/abs/2403.13927}{arXiv:2403.13927}.

  \item[\textup{[Shravan26]}]\label{ref:7293-shravan26}
  S. Shravan, M. Raza, and A. Shlosberg, "Efficient simulation of noisy IQP circuits with amplitude-damping noise," arXiv preprint (2026), version 2, 22 April 2026. \href{https://arxiv.org/abs/2604.05036}{arXiv:2604.05036}.
\end{enumerate}

\subsection*{Comment}

This would give a rigorous route from dissipative loss of conditional correlations to efficient full-distribution sampling at arbitrary depth. Efficient local-observable estimation, decay of ordinary two-point correlations, and a quantum-state conditional-mutual-information bound are not interchangeable with the classical inequality asked here.

\subsection*{Field}

Quantum algorithm

\subsection*{Topic}

Random circuit sampling; Computational complexity and computability

\subsection*{ID}

\texttt{op\_72937f30974953ba}
